%% ma: L12-B08-0074
%% lop: 12
%% chuong: 2
%% bai: 8
%% dang: 12.8.27
%% loai: TLN
%% muc: TH
%% dap_an: "4"
%% nguon: "LaTeX giáo viên gửi 21/10/2026 (M05, Vectơ và hệ trục tọa độ trong không gian, 107 câu), Phần III Câu 14"
%% kien_thuc: "Bài 8 (tọa độ trọng tâm tam giác)"
%% trang_thai: chinh-thuc
%% ghi_chu: "Giáo viên đã duyệt 21/10/2026."
\begin{ex}
Trong không gian $Oxyz$, cho tam giác $MNP$ có $M(2;1;3)$, $N(0;-1;5)$ và $P(4;3;-2)$. Gọi $G(x_G;y_G;z_G)$ là trọng tâm của tam giác $MNP$. Tính $T=x_G\cdot y_G\cdot z_G$.
\shortans{4}
\loigiai{$G\left(\dfrac{2+0+4}3;\dfrac{1-1+3}3;\dfrac{3+5-2}3\right)=(2;1;2)$, $T=4$.}
\end{ex}
